% =====================================================================
% paper/thesis_ch5/sections/related_work.tex
% §5.2 研究背景与定位（Background & Related Work）
% =====================================================================
% rev.2（2026-06-07）：§5.1/§5.2 推倒重写（方案甲；用户拍板）。
%   职责定位：§5.2 = 生物借流 landscape + RL 复现谱系 + 算法背景 + gap 的
%   唯一充分展开处（§5.1 仅一句生物起兴，谱系与 gap 精确区分全在此节）。
%   四段：
%     ¶1 生物借流 landscape：能量场/信息场二分，鱼类充分展开（涡街/尾流/群体涡相位
%        + 侧线"空间分布式"阵列，与 §5.1 一致、反衬 AUV 单点），鸟·浮游一笔印证；
%     ¶2 RL 复现谱系：奠基 → 局部感知导航 → 具身鱼形机器人 → 算法评估；
%     ¶3 算法背景：offline 行为约束谱系（弱→强）+ POMDP 训练-部署不对称/特权信息；
%     ¶4 gap 三支精确区分 + 本章坐标 + forward pointer。
%   红线：深度算法对比（TD3+BC α / ReBRAC dual-penalty / FQL flow-matching）下放
%   §5.6/§5.8 method，此处描述性、不裸写算法专名、不展开公式；gap 第三支收敛
%   "训练与部署阶段均不借特权"（避 asymmetric-AC 稻草人，对齐 §0.1）；两条主线复述
%   不在本节重列，留 §5.1 ¶6；零跨章自足，具体平台/传感硬件留 §5.3 setup。
% rev.3（2026-06-07）：语言层润色（nature-skills 一轮，结构/论点位置/cite key 不动）——
%   ¶2 末句拆 qiu2022（对称性影响策略）/ mecanna2025（多方法系统评估），消轻度类别错配；
%   去口语/翻译腔（减少集体游动能耗 / 借助涡结构 / 即可 / 尚须 / 锚定于）；侧线句分号分层；
%   ¶4 鲜少·至今·删冗余回指。nature-citation 核转述全忠实 + nature-reviewer 把关（合格放行）。
% 前序沿革：rev.1 见 git 历史（与旧 §5.1 hook 一次落地，已被 rev.6/rev.2 整体重写取代）。
% =====================================================================

\section{研究背景与定位}\label{sec:ch5_related}

生物对流场的利用可从两个互补的角度理解：对运动而言，流场既是可资借用的能量来源，也是可供感知的信息来源。就能量而言，鱼类能在涡街与障碍物尾流中降低单体游动的代谢消耗，也能在群体中借尾拍与同伴尾涡之间的相位匹配，减少集体游动的能耗~\citep{liao2022efficiency,disanto2025smarter,li2020vortexphase}；这些现象的共同点在于，高效运动来自身体运动与流动结构在时空上的耦合，而非单纯增大推力。就信息而言，支撑上述行为的鱼类侧线是一种沿体表空间分布的感受阵列，而非单点的速度计~\citep{coombs2020rheotaxis}；这一结构提示，传感在空间上的组织方式本身，可能决定借流行为的可行性。同一原则在更广的类群与尺度上亦有印证：鸟类借热上升气流实现长距离滑翔，浮游生物则响应湍流中的局部梯度完成输运~\citep{flato2024soaring,mousavi2024zooplankton}。

强化学习把这些借流现象从生物观察转化为可建模、可学习的控制问题。奠基性的工作让智能微游泳体在解析流场中学习游动方向以避开不利流区，并借助涡结构提升集体游动的推进效率~\citep{colabrese2017flownav,verma2018pnas}。其后研究的重心转向局部感知：固定速度的游泳体仅凭局部流速，即可在非定常涡流中学到接近全局最优的高效轨迹~\citep{gunnarson2021ncomm}；而在自我中心坐标系下，单一的局部流速往往不足以支撑可靠学习，尚须辅以局部的流场梯度；这一结论恰与鱼类侧线的空间阵列结构相呼应~\citep{jiao2025gradients}。更贴近工程实现的具身研究让柔性鱼形机器人在涡街与圆柱尾流中学习稳健的导航与姿态控制~\citep{zhu2022fishlike,feng2025roboticfish}；相关研究还表明，流场对称性等物理结构会影响所学策略~\citep{qiu2022symmetries}，并对复杂、局部可观测流场中不同方法的表现作了系统评估~\citep{mecanna2025assessment}。

本章所用的方法工具取自两条相对独立的研究脉络。离线强化学习提供了从固定数据中稳健学习的手段：为遏制分布外动作引起的价值高估，行为约束类方法以不同强度将学得的策略锚定于数据分布，从对动作的显式正则，到以生成式模型刻画行为分布的隐式先验，构成一条由弱到强的谱系~\citep{levine2020offline,fujimoto2021td3bc,tarasov2023rebrac,park2025fql}。部分可观测下的训练与部署不对称则提供了在训练期借助额外信息的手段：当训练阶段可获得超出部署传感的额外观测，即特权信息时，非对称 actor-critic 允许评价器利用这些信息以稳定学习，而策略本身始终只依赖部署阶段可得的输入~\citep{pinto2017asymmetric}。这两条脉络分别解决了从固定数据中学习、在训练期利用额外信息这两个问题，却都未针对结构化流场中的单点物理传感约束作系统考察。

将上述脉络并置，便显出一处共同的空白。借流强化学习的研究大多依赖多于单点的流场观测，或允许智能体在线反复试错；离线强化学习的研究多在标准基准上验证，鲜少触及单点局部可观测这一物理传感约束；利用特权信息的工作则通常在训练阶段为评价器引入超出部署传感的额外观测。三者交叉之处，即仅有单点局部感知、只能纯离线学习、且训练与部署阶段均不借助任何特权信息的情形，恰是受约束平台的真实处境，却至今缺乏系统刻画。本章正定位于这一交叉处：不提出新的学习算法，而是在统一的受控设定与多随机种子的统计检验下，辨明局部感知约束下可行性的来源及其失效边界，并揭示离线学习内部决定算法高下的真正因素。这些发现共同指向的机制，将由后续各节逐一展开。
